\subsection{Definición estructural y función de la constante gravitatoria}
\label{sec:ii-definicion-gravedad}

La realización común recibe primero la longitud
\(L=(5759/23040)\alpha_{\HMT}^{16}L_*\) de
\eqref{g26:eq:longitud-unica-es}. Su duración elemental es
\(t_0=\pi_{\HMT}L/(54c_{\rm int})\); por tanto,
\(R_{108}=54c_{\rm int}t_0/\pi_{\HMT}=L\).
La igualdad con la longitud de Planck de esta realización,
\(\ell_P=L\), se obtiene después de determinar el acoplamiento.
La referencia positiva \(L_*\) se mantiene explícita; la evaluación
en metros utiliza la sección declarada \(L_*=1\,\mathrm m\).

\begin{definition}[Sección gravitatoria contragrediente de la acción]
Para una sección positiva \(\hbar_a\), una velocidad
\(c_{\rm int}>0\) y el radio cronológico \(R_{108}\) de la
sección~\ref{sec:gravity}, la sección gravitatoria circular es
\begin{equation}
 G_a=\frac{c_{\rm int}^{3}R_{108}^{2}}{\hbar_a},
 \qquad R_{108}=\frac{54}{\pi}\ell_0,
 \qquad \ell_0=c_{\rm int}t_0.
 \label{eq:ii-definicion-gravedad-seccion}
\end{equation}
Es la sección recíproca de la acción que realiza la coincidencia
entre el radio gravitatorio reducido del ciclo y su radio
cronológico, en la convención \(r_g=Gm/c_{\rm int}^2\).
\end{definition}

La proposición de coincidencia de radios de la
sección~\ref{sec:gravity} expresa su unicidad en la colección
\(G_a(\vartheta)=\vartheta c_{\rm int}^{5}t_0^2/\hbar_a\).
La condición circular fija \(\vartheta=(54/\pi)^2\) y conduce
a \eqref{eq:ii-definicion-gravedad-seccion}. Esa condición y las
secciones dimensionales positivas forman parte de la realización.
La prueba operatoria de \ref{g26:sec:rigidez-radial-es} amplía esa
expresión al carácter positivo completo: R4 define la composición
longitudinal y R5 su métrica radial; R8 identifica el operador resultante
con el radio gravitatorio reducido. Con esos datos, el mismo coeficiente
\(G_{\rm ret}=c_{\rm int}^3L^2/\hbar_{\rm ret}\) actúa en todos
los modos positivos admisibles. La fórmula circular es su especialización
al carácter identidad y al reloj conjugado de la longitud.
Manteniéndolas comunes entre las dos orientaciones, se obtiene
\begin{equation}
 \frac{G_{\mathrm{pre}}}{G_{\mathrm{ret}}}
 =R_{\mathrm{act}}^{-1},
 \qquad
 \frac{G_a\hbar_a}{c_{\rm int}^{3}}
 =R_{108}^{2}=L^2=\ell_P^2.
 \label{eq:ii-definicion-gravedad-area}
\end{equation}
La prueba es la sustitución directa de
\eqref{eq:ii-definicion-gravedad-seccion}; el cambio de acción
se compensa por el cambio recíproco de \(G_a\).

Esta estructura asigna a la gravitación una función precisa:
transportar la sección de acción a una unidad geométrica común.
En \eqref{eq:holst}, \(G\) fija la normalización del acoplamiento;
en el operador de área, \(G\hbar/c^3\) fija la unidad areal.
La reciprocidad conserva esa unidad mientras cambia la lectura
orientada de acción. El estado reducido y sus ocupaciones aportan
después la información y el entrelazamiento asociados a la frontera.

El contenido tensorial del acoplamiento se conserva en
\(\mathbb G_{5,a}=G_aP^{\mathrm{ico}}_5\), construido en
\ref{sec:ii-pentadica}. La sección \(G_a\) determina su magnitud común;
el proyector selecciona el portador pentacomponente y \(\Gamma_5\)
transporta su orientación al espacio de tensores simétricos sin traza.
La base de \eqref{ii:pent:eq:gravity-helicity2-real}--\eqref{ii:pent:eq:gravity-helicity0-real}
explicita las cinco componentes. La proyección \(\Lambda(n)\) define
después la reducción transversal. Magnitud, portador y reducción son
operaciones relacionadas, con dominios diferentes dentro de esta realización.
